% !TeX root = ../cosmology_summary.tex

%% --------------------------------------------------------
\chapter{Хаотична інфляція Лінде}
\label{sec:chaotic}
%% --------------------------------------------------------

Лінде~\cite{linde_1983} запропонував радикально простий погляд
на початкові умови: значення $\phi$ у різних точках раннього
Всесвіту є хаотичним~--- немає потреби в особливому «стартовому»
станні. Найпростіший потенціал:
\begin{equation*}
  V(\phi) = \tfrac{1}{2}m^2\phi^2.
\end{equation*}

Де $\phi$ велике і $\varepsilon \ll 1$~--- там інфляція іде довго
і виникає величезна причинно зв'язана «бульбашка». Де $\phi$ мале~---
інфляція швидко завершується. Наш спостережуваний Всесвіт~---
одна з таких бульбашок.

\textbf{Структура всередині однієї бульбашки.} Коли інфляція
в нашій бульбашці завершилась, простір всередині неї вже був
практично однорідним~--- розширення загладило всі початкові
неоднорідності $\phi$. Але «практично» не означає «абсолютно»:
залишились крихітні квантові флуктуації, заморожені під час
інфляції. Саме вони стали насінням галактик і є тим, що ми
бачимо на карті CMB.

Всередині однієї бульбашки існує єдиний причинно зв'язаний
простір-час~--- один Всесвіт. Але його розмір після $N \sim 60$
$e$-folds становить $\sim e^{60}$ разів більше за наш
спостережуваний горизонт ($\sim 46$~Гпк). Тому навіть у
межах «нашої бульбашки» існують області, з якими ми ніколи
не матимемо причинного контакту~--- вони назавжди за
нашим горизонтом подій.

\textbf{Між бульбашками.} Простір між різними бульбашками
продовжує розширюватись інфляційно~--- швидше за світло.
Дві бульбашки, що утворились незалежно, розділені областю
вічної інфляції і принципово не можуть обмінятись сигналом.
Це не просто «дуже далеко»~--- між ними немає спільного
майбутнього світлового конуса. У строгому сенсі загальної теорії відносності вони
належать різним причинно відокремленим регіонам
одного глобального простору-часу, а не різним «всесвітам»~---
хоча в популярній літературі їх так і називають.

\textbf{Вічна інфляція.} Квантові флуктуації інфлатона на
суперхабблових масштабах можуть локально \emph{збільшити} $\phi$,
повернувши його назад на схил, якщо флуктуація достатньо велика.
Класичний зсув поля за один $e$-fold: $|\delta\phi_{\rm class}|
\sim |\dot\phi|/H \approx \sqrt{2\varepsilon}\,M_{\rm Pl}$.
Квантова флуктуація за той самий час: $\delta\phi_{\rm quant}
\sim H/(2\pi)$. Умова «самовідтворення»~---
квантова флуктуація перевищує класичне зміщення:
\begin{equation*}
  \frac{H}{2\pi} \gtrsim \sqrt{2\varepsilon}\,M_{\rm Pl}
  \quad\Leftrightarrow\quad
  H^2 \gtrsim M_{\rm Pl}^2
  \quad\Leftrightarrow\quad
  V(\phi) \gtrsim M_{\rm Pl}^4,
\end{equation*}
де в останньому переході використано $H^2 \approx V/(3M_{\rm Pl}^2)$
і $\varepsilon \lesssim 1$.
У таких областях інфляція ніколи не зупиняється глобально:
поки одна бульбашка завершує інфляцію і термалізується,
сусідні регіони продовжують розширюватись, породжуючи нові
бульбашки. Це \textbf{вічна інфляція}~--- мультивсесвіт
як математичний наслідок, а не постулат.

Строгий теоретичний опис цього динамічного режиму забезпечується стохастичним
підходом Старобінського. У цій формалізації квантові флуктуації короткохвильових
мод розглядаються як випадкове джерело (білий шум), що діє на класичну
довгохвильову координату поля. Еволюція інфлатона набуває вигляду стохастичного
рівняння Ланжевена, де класичний дрейф по схилу потенціалу конкурує з квантовою
дифузією. Коли поле перебуває поблизу планківської межі, дифузійний внесок
повністю домінує, що призводить до квантового блукання поля вгору по потенціалу
і робить процес генерації нових областей простору просторово нескінченним.

\textbf{Статус.} Потенціал $\tfrac{1}{2}m^2\phi^2$ передбачає
$r \approx 8/N \approx 0.13$~--- вище поточної межі BICEP/Keck
(табл.~\ref{tab:inflation_models}), тому ця конкретна модель
виключена. Але ідея хаотичних початкових умов і вічної інфляції
залишається живою у більш пологих потенціалах. Сучасним переосмисленням
парадигми Лінде є клас моделей $\alpha$-аттракторів та інфляція з немінімальним
зв'язком інфлатона з кривиною простіру-часу. У цих сценаріях при великих
значеннях поля потенціал природно виходить на асимптотичне плато, що дозволяє
зберегти концепцію хаотичного старту високої енергії, але радикально знижує
амплітуду тензорних збуджень $r$, повністю узгоджуючи теорію з останніми даними
Planck та BICEP/Keck.

Разом з тим жодна з «успішних» моделей~--- тих, що проходять
обмеження Planck на $n_s$ і $r$~--- не передбачає нічого
особливого на малих масштабах $k \gtrsim 1\,h/$Мпк. Саме там
первинний спектр залишається практично необмеженим
спостереженнями. І саме там у 2022--2023 роках телескоп JWST
зафіксував надлишок масивних галактик при $z \sim 10$--$12$:
об'єктів, які за стандартним плоским спектром просто не встигли
б вирости за відведений час~\cite{finkelstein_2023,
boylan_kolchin_2023}.

Природне пояснення~--- локальне підсилення первинного спектра
саме на цих масштабах. Якщо потенціал інфлатона має особливість
(сходинку або ділянку ультра-повільного кочення) у потрібному
місці, флуктуації, що виходять за горизонт в цей момент,
зазнають різкого підсилення~--- формується \emph{зламаний спектр
потужності} (рис.~\ref{fig:jwst_feature}). Чи здатна ця ідея витримати зіставлення з повним
масивом даних Planck і JWST~--- саме це є предметом
наступного розділу.

%% --------------------------------------------------------
%% РИСУНОК: МОДИФІКАЦІЯ ПОТЕНЦІАЛУ ТА ЗЛАМАНИЙ СПЕКТР
%% --------------------------------------------------------
\begin{figure}[htbp]
  \centering
  \begin{tikzpicture}
    \pgfplotsset{compat=1.18}

    % Лівий графік: Потенціал з особливістю (USR)
    \begin{axis}[
        width=0.48\textwidth,
        height=5.5cm,
        axis lines=left,
        xlabel={$\phi$},
        ylabel={$V(\phi)$},
        ymin=0, ymax=4.0,
        xmin=0, xmax=4.5,
        ticks=none,
        title={Потенціал інфлатона}
      ]
      % Монотонна функція з вираженим плато (похідна в районі x=2 падає майже до нуля)
      \addplot[thick, color=blue, domain=0:4.2, samples=100, smooth]
        {0.1*x^3 - 0.6*x^2 + 1.3*x + 0.5};

      % Позначення ділянки особливості (USR)
      \draw[dashed, gray] (axis cs:1.3,0) -- (axis cs:1.3,1.5);
      \draw[dashed, gray] (axis cs:2.7,0) -- (axis cs:2.7,1.5);
      \draw[<->, >=stealth, gray] (axis cs:1.3,0.4) -- node[above, font=\scriptsize, text=black] {USR} (axis cs:2.7,0.4);

      % Акуратна коротка стрілка-вказівка
      \node[above, font=\scriptsize] at (axis cs:1.6,2.7) {Особливість (сходинка)};
      \draw[->, >=stealth, thick, darkgray] (axis cs:1.6,2.6) -- (axis cs:2.0,1.55);
    \end{axis}

    % Правий графік: Зламаний спектр потужності
    \begin{axis}[
        width=0.48\textwidth,
        height=5.5cm,
        xshift=6.8cm,
        axis lines=left,
        xlabel={$\ln k$},
        ylabel={$\ln \mathcal{P}_\zeta(k)$},
        ymin=-1, ymax=4.5,
        xmin=0, xmax=5,
        ticks=none,
        title={Первинний спектр потужності}
      ]
      % Масштаби СМВ (плато) та різкий пік на малих масштабах
      \addplot[thick, color=red, domain=0:5, samples=150, smooth]
        {0.5 + 3.2*exp(-(x-3.4)^2/0.18)};

      % Позначення зон обмежень
      \draw[dashed, gray] (axis cs:0,0.5) -- (axis cs:2.2,0.5);
      \node[below, font=\scriptsize, text=darkgray] at (axis cs:1.1,0.5) {CMB (Planck)};

      \draw[dashed, gray] (axis cs:3.4,-1) -- (axis cs:3.4,3.7);
      \node[above, font=\scriptsize, text=red] at (axis cs:3.4,3.7) {Пік JWST};
    \end{axis}
  \end{tikzpicture}
  \caption{Схематичне зображення механізму генерації зламаного спектра.
           Ліворуч: потенціал інфлатона з короткою ділянкою ультра-повільного кочення (USR),
           де класична швидкість поля падає. Праворуч: відповідний спектр потужності кривини
           $\mathcal{P}_\zeta(k)$, що зберігає масштабну інваріантність на масштабах СМВ,
           але має локальний надлишок амплітуди на високих хвильових числах.}
  \label{fig:jwst_feature}
\end{figure}